\section{Introduction}
\label{sec:intro}

Unsupervised text clustering serves as a cornerstone for exploring unlabelled textual data.
However, traditional algorithmic approaches operate under a rigid paradigm, assuming that a singular, objectively correct group assignment exists \cite{clusterllm, dial_in_llm}.
In practical applications, text categorization is highly subjective; 
two users analyzing the same dataset may demand completely divergent cluster structures based on their immediate task constraints, domain expertise, or semantic perspectives \cite{eliciting_human_preferences}. 
Consequently, static clustering solutions often fail to align with idiosyncratic human intents.

To address this gap, recent paradigms have shifted toward integrating human-in-the-loop guidance within clustering workflows.
Concurrently, Large Language Models (LLMs) have emerged as powerful conduits for natural language preference elicitation, uncovering latent user nuances far more dynamically than passive feedback mechanics \cite{eliciting_human_preferences, to_gate, asking_clarifying_questions}. 
By positioning an LLM as an interviewing agent, a system can iteratively converse with a user, capture their architectural preferences for the data, and adapt downstream clustering behaviors accordingly.

This project introduces \textit{Cluster Around and Find Out}, a conversational text clustering system designed to bridge user intent and document representation. 
While clustering remains central to our pipeline, our empirical evaluation targets the conversational interface itself rather than raw cluster geometry against artificial baselines. 
Specifically, we investigate two core research questions:
\begin{itemize}
    \item \textbf{Q1:} Which questioning paradigm (yes/no, multiple-choice, or open-ended) yields the most efficient path to convergence while minimizing human cognitive load?
    \item \textbf{Q2:} Does incorporating task-specific instructions derived from conversational preferences materially improve the subjective faithfulness of the resulting clusters compared to an uninstructed baseline?
\end{itemize}

We formalize our investigation through a rigorous study design utilizing a simulated user with 3 distinct user personas and 6 specialized clustering goals across a subset of the StackOverflow dataset. 
We explicitly test the hypothesis that multiple-choice questions strike the premier equilibrium—optimizing the volume of preference data harvested per turn without placing excessive linguistic or cognitive demands on the user. 

The remainder of this report is organized as follows: 
Section~\ref{sec:related_work} contextualizes our approach within the state-of-the-art literature; 
Section~\ref{sec:study_design} and Section~\ref{sec:methodology} outline our experimental variables and system architecture; 
Section~\ref{sec:results} and Section~\ref{sec:discussion} present our bootstrap-tested findings; 
and Section~\ref{sec:limitations} addresses key confounding factors and future avenues.